% Rol 2. Requiere amsmath; insertar con % Rol 2. Requiere amsmath; insertar con % Rol 2. Requiere amsmath; insertar con % Rol 2. Requiere amsmath; insertar con \input{seccion_6_diseno_experimental}.
\section{Diseño experimental y evaluación}
La unidad de evaluación es una predicción a una hora; persistencia y modelos usan las mismas filas completas. Se conserva la división cronológica de la sección 4. Dentro del entrenamiento se ajusta $\alpha$ con tres ventanas de origen expansivo: validación interna de julio--septiembre de 2022, octubre--diciembre de 2022 y enero--marzo de 2023. Cada ventana se entrena con las observaciones anteriores a ella. Los cortes se aplican por fecha de la respuesta, incluso si hay horas faltantes. Se prueba $\alpha\in\{0{,}001,0{,}01,0{,}1,1,10,100,1000\}$, por separado para cada modelo y conjunto de variables. Se minimiza el promedio de MAE de las tres ventanas. El escalador se ajusta por separado en cada ventana dentro del pipeline.

Tras la búsqueda, cada variante se ajusta con todo el entrenamiento y se evalúa en validación (julio--diciembre de 2023). Se compara con persistencia sobre las mismas observaciones y se selecciona el modelo lineal con menor MAE de validación. Se evalúan ese modelo y persistencia en prueba (enero--abril de 2024), sin reajustar con validación. La prueba no se emplea para seleccionar hiperparámetros ni variantes. La comparación con y sin calendario cambia solamente ese conjunto de variables, con búsqueda independiente. Los resultados y predicciones se registran en archivos CSV.

Sean $y_i$ las observaciones, $\hat y_i$ los pronósticos y $\bar y$ la media observada del bloque. Se reportan MAE $=n^{-1}\sum |y_i-\hat y_i|$, MSE $=n^{-1}\sum(y_i-\hat y_i)^2$, RMSE $=\sqrt{\mathrm{MSE}}$ y $R^2=1-\sum(y_i-\hat y_i)^2/\sum(y_i-\bar y)^2$. El MAE, en las unidades de concentración, es el criterio principal por su interpretación directa; RMSE penaliza más los errores grandes y conserva las unidades originales. MSE expresa esa penalización en unidades cuadradas y $R^2$ compara con la media observada del bloque, no con persistencia; puede ser negativo. Para MAE, RMSE y MSE un valor menor es mejor, y para $R^2$ uno mayor es mejor. La mejora porcentual frente a persistencia se calcula como $100(1-\mathrm{MAE}_{modelo}/\mathrm{MAE}_{persistencia})$ en el mismo bloque.

El EDA previo examinó todo el periodo, incluida prueba: aunque los cortes se fijaron antes de entrenar los modelos, esto limita la independencia estricta de la evaluación final. La exclusión de ventanas incompletas también limita la generalización a periodos con mediciones disponibles. Las diferencias entre periodos deben interpretarse teniendo en cuenta el cambio de distribución de concentraciones.
.
\section{Diseño experimental y evaluación}
La unidad de evaluación es una predicción a una hora; persistencia y modelos usan las mismas filas completas. Se conserva la división cronológica de la sección 4. Dentro del entrenamiento se ajusta $\alpha$ con tres ventanas de origen expansivo: validación interna de julio--septiembre de 2022, octubre--diciembre de 2022 y enero--marzo de 2023. Cada ventana se entrena con las observaciones anteriores a ella. Los cortes se aplican por fecha de la respuesta, incluso si hay horas faltantes. Se prueba $\alpha\in\{0{,}001,0{,}01,0{,}1,1,10,100,1000\}$, por separado para cada modelo y conjunto de variables. Se minimiza el promedio de MAE de las tres ventanas. El escalador se ajusta por separado en cada ventana dentro del pipeline.

Tras la búsqueda, cada variante se ajusta con todo el entrenamiento y se evalúa en validación (julio--diciembre de 2023). Se compara con persistencia sobre las mismas observaciones y se selecciona el modelo lineal con menor MAE de validación. Se evalúan ese modelo y persistencia en prueba (enero--abril de 2024), sin reajustar con validación. La prueba no se emplea para seleccionar hiperparámetros ni variantes. La comparación con y sin calendario cambia solamente ese conjunto de variables, con búsqueda independiente. Los resultados y predicciones se registran en archivos CSV.

Sean $y_i$ las observaciones, $\hat y_i$ los pronósticos y $\bar y$ la media observada del bloque. Se reportan MAE $=n^{-1}\sum |y_i-\hat y_i|$, MSE $=n^{-1}\sum(y_i-\hat y_i)^2$, RMSE $=\sqrt{\mathrm{MSE}}$ y $R^2=1-\sum(y_i-\hat y_i)^2/\sum(y_i-\bar y)^2$. El MAE, en las unidades de concentración, es el criterio principal por su interpretación directa; RMSE penaliza más los errores grandes y conserva las unidades originales. MSE expresa esa penalización en unidades cuadradas y $R^2$ compara con la media observada del bloque, no con persistencia; puede ser negativo. Para MAE, RMSE y MSE un valor menor es mejor, y para $R^2$ uno mayor es mejor. La mejora porcentual frente a persistencia se calcula como $100(1-\mathrm{MAE}_{modelo}/\mathrm{MAE}_{persistencia})$ en el mismo bloque.

El EDA previo examinó todo el periodo, incluida prueba: aunque los cortes se fijaron antes de entrenar los modelos, esto limita la independencia estricta de la evaluación final. La exclusión de ventanas incompletas también limita la generalización a periodos con mediciones disponibles. Las diferencias entre periodos deben interpretarse teniendo en cuenta el cambio de distribución de concentraciones.
.
\section{Diseño experimental y evaluación}
La unidad de evaluación es una predicción a una hora; persistencia y modelos usan las mismas filas completas. Se conserva la división cronológica de la sección 4. Dentro del entrenamiento se ajusta $\alpha$ con tres ventanas de origen expansivo: validación interna de julio--septiembre de 2022, octubre--diciembre de 2022 y enero--marzo de 2023. Cada ventana se entrena con las observaciones anteriores a ella. Los cortes se aplican por fecha de la respuesta, incluso si hay horas faltantes. Se prueba $\alpha\in\{0{,}001,0{,}01,0{,}1,1,10,100,1000\}$, por separado para cada modelo y conjunto de variables. Se minimiza el promedio de MAE de las tres ventanas. El escalador se ajusta por separado en cada ventana dentro del pipeline.

Tras la búsqueda, cada variante se ajusta con todo el entrenamiento y se evalúa en validación (julio--diciembre de 2023). Se compara con persistencia sobre las mismas observaciones y se selecciona el modelo lineal con menor MAE de validación. Se evalúan ese modelo y persistencia en prueba (enero--abril de 2024), sin reajustar con validación. La prueba no se emplea para seleccionar hiperparámetros ni variantes. La comparación con y sin calendario cambia solamente ese conjunto de variables, con búsqueda independiente. Los resultados y predicciones se registran en archivos CSV.

Sean $y_i$ las observaciones, $\hat y_i$ los pronósticos y $\bar y$ la media observada del bloque. Se reportan MAE $=n^{-1}\sum |y_i-\hat y_i|$, MSE $=n^{-1}\sum(y_i-\hat y_i)^2$, RMSE $=\sqrt{\mathrm{MSE}}$ y $R^2=1-\sum(y_i-\hat y_i)^2/\sum(y_i-\bar y)^2$. El MAE, en las unidades de concentración, es el criterio principal por su interpretación directa; RMSE penaliza más los errores grandes y conserva las unidades originales. MSE expresa esa penalización en unidades cuadradas y $R^2$ compara con la media observada del bloque, no con persistencia; puede ser negativo. Para MAE, RMSE y MSE un valor menor es mejor, y para $R^2$ uno mayor es mejor. La mejora porcentual frente a persistencia se calcula como $100(1-\mathrm{MAE}_{modelo}/\mathrm{MAE}_{persistencia})$ en el mismo bloque.

El EDA previo examinó todo el periodo, incluida prueba: aunque los cortes se fijaron antes de entrenar los modelos, esto limita la independencia estricta de la evaluación final. La exclusión de ventanas incompletas también limita la generalización a periodos con mediciones disponibles. Las diferencias entre periodos deben interpretarse teniendo en cuenta el cambio de distribución de concentraciones.
.
\section{Diseño experimental y evaluación}
La unidad de evaluación es una predicción a una hora; persistencia y modelos usan las mismas filas completas. Se conserva la división cronológica de la sección 4. Dentro del entrenamiento se ajusta $\alpha$ con tres ventanas de origen expansivo: validación interna de julio--septiembre de 2022, octubre--diciembre de 2022 y enero--marzo de 2023. Cada ventana se entrena con las observaciones anteriores a ella. Los cortes se aplican por fecha de la respuesta, incluso si hay horas faltantes. Se prueba $\alpha\in\{0{,}001,0{,}01,0{,}1,1,10,100,1000\}$, por separado para cada modelo y conjunto de variables. Se minimiza el promedio de MAE de las tres ventanas. El escalador se ajusta por separado en cada ventana dentro del pipeline.

Tras la búsqueda, cada variante se ajusta con todo el entrenamiento y se evalúa en validación (julio--diciembre de 2023). Se compara con persistencia sobre las mismas observaciones y se selecciona el modelo lineal con menor MAE de validación. Se evalúan ese modelo y persistencia en prueba (enero--abril de 2024), sin reajustar con validación. La prueba no se emplea para seleccionar hiperparámetros ni variantes. La comparación con y sin calendario cambia solamente ese conjunto de variables, con búsqueda independiente. Los resultados y predicciones se registran en archivos CSV.

Sean $y_i$ las observaciones, $\hat y_i$ los pronósticos y $\bar y$ la media observada del bloque. Se reportan MAE $=n^{-1}\sum |y_i-\hat y_i|$, MSE $=n^{-1}\sum(y_i-\hat y_i)^2$, RMSE $=\sqrt{\mathrm{MSE}}$ y $R^2=1-\sum(y_i-\hat y_i)^2/\sum(y_i-\bar y)^2$. El MAE, en las unidades de concentración, es el criterio principal por su interpretación directa; RMSE penaliza más los errores grandes y conserva las unidades originales. MSE expresa esa penalización en unidades cuadradas y $R^2$ compara con la media observada del bloque, no con persistencia; puede ser negativo. Para MAE, RMSE y MSE un valor menor es mejor, y para $R^2$ uno mayor es mejor. La mejora porcentual frente a persistencia se calcula como $100(1-\mathrm{MAE}_{modelo}/\mathrm{MAE}_{persistencia})$ en el mismo bloque.

El EDA previo examinó todo el periodo, incluida prueba: aunque los cortes se fijaron antes de entrenar los modelos, esto limita la independencia estricta de la evaluación final. La exclusión de ventanas incompletas también limita la generalización a periodos con mediciones disponibles. Las diferencias entre periodos deben interpretarse teniendo en cuenta el cambio de distribución de concentraciones.
